\chapter{Bronchoalveolar Lavage}

\section*{Overview}
Bronchoalveolar lavage (BAL) is a procedure performed during bronchoscopy in which sterile saline is instilled into a segment of the lung and then suctioned back out. This washes the alveoli and collects cells, microorganisms, and other material for laboratory analysis.

\section*{Alternative Names}
BAL, bronchoalveolar lavage, lung lavage, bronchoscopic lavage


\section*{Principle}
BAL fluid is analyzed for total cell count and differential (macrophages, lymphocytes, neutrophils, eosinophils) by cytocentrifugation and Wright-Giemsa staining. Flow cytometry subsets lymphocytes. Microbiology includes Gram stain, fungal and AFB smears, cultures, and PCR. Cytology screens for malignant cells.

\section*{History}
Bronchoalveolar lavage was developed in the 1970s with the advent of flexible bronchoscopy. It enabled sampling of the lower respiratory tract for diagnosing infections, interstitial lung disease, and malignancy.

\section*{Why It Is Done}
Your doctor orders this test when they need to examine the cells and organisms deep within your lungs. It is performed during a bronchoscopy, where a thin camera is passed through your mouth or nose into your lungs. It answers the question: what is causing your lung disease --- an infection, an autoimmune condition, or an occupational exposure?

\section*{Is This a Primary or Secondary Test or Procedure}
Secondary. It is performed when less invasive tests like sputum samples or imaging have not provided a clear answer.

\section*{How It Is Performed}
During bronchoscopy, a thin flexible tube with a camera is passed through your mouth or nose, down your throat, and into your lungs. Saline is injected into a small section of your lung and then suctioned back out. This washes the alveoli (tiny air sacs) and collects cells, organisms, and other material for analysis. You receive sedation during the procedure, so you should not feel pain.

\section*{Preparation}
Fasting and sedation instructions come from the bronchoscopy team. Arrange transportation if sedation is used. Tell the team about anticoagulants, antiplatelet medicines, and bleeding disorders; do not stop blood-thinning medicines unless the prescribing and procedural teams provide a coordinated plan.

\section*{Contraindications and Precautions}
This procedure has the following contraindications and precautions:

\begin{itemize}
    \item \textbf{Relative contraindications}: Severe hypoxemia (PaO2 below 60 mmHg or significant oxygen requirement), unstable asthma with active bronchospasm, uncontrolled coagulopathy or severe thrombocytopenia, severe pulmonary hypertension, and hemodynamic instability.
    \item \textbf{Precautions}: Correct coagulopathy when feasible before the procedure; use continuous pulse oximetry and supplemental oxygen during and after the procedure; limit the volume of saline instilled in patients with marginal respiratory reserve; coordinate any anticoagulant plan with the prescribing and procedural teams.
    \item \textbf{Special populations}: Children require smaller lavage volumes and shorter procedure time; pregnant patients should undergo the procedure only when the diagnostic benefit outweighs fetal risk.
\end{itemize}

\section*{Risks}
Transient hypoxemia, localized fever, minor bleeding, or bronchospasm.

\section*{Aftercare and Recovery}
Vital signs are monitored for 1-2 hours; avoid eating or drinking until your gag reflex returns.

\section*{Outcome and Follow-up}
The lavage fluid is analyzed for cell count and differential, microbiology, and cytology. A lymphocytic pattern (often with a low CD4:CD8 ratio) suggests sarcoidosis or hypersensitivity pneumonitis, a neutrophilic predominance suggests bacterial infection, and an eosinophilic pattern suggests eosinophilic lung disease or drug reaction. Positive cultures or PCR identify the causative organism.

\section*{Expected Results}
Normal: No predominant pathogen. Mixed respiratory flora. WBC differential varies by lavage site. No malignant cells.

\begin{center}
\begin{tabular}{ll}
\toprule
\textbf{Population} & \textbf{Reference Range} \\
\midrule
Adult Male & No predominant pathogen; mixed flora; no malignant cells \\
Adult Female & Same as male \\
Children & Same as adult \\
Elderly & Same as adult \\
Pregnancy & Same as non-pregnant \\
\bottomrule
\end{tabular}
\end{center}

\section*{Follow-up}
Follow-up depends on the findings:
\begin{itemize}
    \item \textbf{Antimicrobial therapy:} When infection is identified, targeted treatment is guided by culture and sensitivity results.
    \item \textbf{Repeat bronchoscopy:} May be repeated if the initial lavage is non-diagnostic or if symptoms persist.
    \item \textbf{Surgical lung biopsy:} Considered when the BAL remains non-diagnostic and interstitial lung disease is suspected.
    \item \textbf{Specialist consultation:} Pulmonology referral for integrated interpretation and management planning.
\end{itemize}
\section*{When to Seek Medical Care}
Follow the procedure team's specific instructions. Seek urgent medical assessment for severe or worsening pain, uncontrolled bleeding, fever or signs of infection, trouble breathing, chest pain, fainting, new weakness or confusion, inability to urinate, or any rapidly worsening symptom. The appropriate warning signs depend on the procedure and the patient's medical condition.

\section*{References}
\begin{enumerate}
  \item MedlinePlus. Bronchoalveolar Lavage. U.S. National Library of Medicine; 2025.
  \item Current Medical Diagnosis and Treatment. McGraw-Hill; 2025.
\end{enumerate}

\clearpage
